\section{Methodology}
\label{sec:method}

In this section, we present the mathematical formulation of weight-space model merging and task vectors. We then introduce our proposed search space parameterizations and contrast our offline validation optimization with online test-time adaptation.

\subsection{Problem Formulation and Task Vectors}
Let $W_{base} \in \mathbb{R}^D$ represent the weights of a pre-trained base network (e.g., a Vision Transformer). Suppose we are given $K$ expert models, where each expert $k \in \{1, \dots, K\}$ has been fine-tuned from $W_{base}$ to specialize in task $k$. The task vector $V_k \in \mathbb{R}^D$ representing the specialized capabilities of expert $k$ is defined as the parameter difference from the base model \cite{ilharco2022editing}:
\begin{equation}
V_k = W_k - W_{base}
\end{equation}

For a network with $L$ layers, we can partition the base model weights and the task vectors layer-by-layer. Let $W_{base}^{(l)}$ and $V_k^{(l)}$ denote the base weights and the task vector of expert $k$ at layer $l \in \{1, \dots, L\}$, respectively.
Our goal is to construct a single, static merged model $W_{merged}$ that performs well across all $K$ tasks. The merged model weights $W_{merged}^{(l)}$ at layer $l$ are defined as:
\begin{equation}
W_{merged}^{(l)}(\theta) = W_{base}^{(l)} + \sum_{k=1}^K \alpha_k(l; \theta) V_k^{(l)}
\end{equation}
where $\alpha_k(l; \theta) \in \mathbb{R}$ is the merging coefficient for task $k$ at layer $l$, parameterized by $\theta$.

\subsection{Search Space Parameterizations}
We explore three distinct parameterizations $\theta$ for the merging coefficients, representing different trade-offs between expressive capacity and optimization complexity.

\subsubsection{Global Task-Wise Coefficients (GT-Merge)}
To prevent high-dimensional overfitting, we first define a highly constrained, low-dimensional search space where the merging coefficient for task $k$ is constant across all layers. In GT-Merge, we set:
\begin{equation}
\alpha_k(l; \theta) = \alpha_k \quad \forall l \in \{1, \dots, L\}
\end{equation}
The parameter vector is $\theta = [\alpha_1, \dots, \alpha_K]^T \in \mathbb{R}^K$. For a model with $K$ tasks, this search space has a dimensionality of exactly $K$, which is completely independent of the network depth $L$.

\subsubsection{Polynomial Coefficient Profiles (Poly-Val-Merge)}
Following PolyMerge \cite{polymerge}, we model the merging coefficient for task $k$ as a continuous polynomial function of the normalized network depth $l/L$. This parameterization captures the layer-wise sensitivity of deep representations while maintaining a low parameter count. For a polynomial of degree $d$, we define:
\begin{equation}
\alpha_k(l; \theta) = \sum_{j=0}^d c_{kj} \left(\frac{l}{L}\right)^j
\end{equation}
The parameter matrix is $\theta = [c_{kj}] \in \mathbb{R}^{K \times (d+1)}$. The total dimensionality of this search space is $K(d+1)$. When $d=0$, this collapses to GT-Merge. For $d \in \{1, 2, 3\}$, we can model linear, quadratic, and cubic trajectories of merging coefficients across the layers of the model.

\subsubsection{Unconstrained Layer-Wise Search Space}
As a comparison, we also consider the unconstrained layer-wise search space used by online AdaMerging \cite{adamerging}. Here, each layer for each task is assigned an independent scalar coefficient:
\begin{equation}
\alpha_k(l; \theta) = \alpha_{k, l}
\end{equation}
The parameter vector is $\theta = \{\alpha_{k, l}\}_{k=1, l=1}^{K, L} \in \mathbb{R}^{K \times L}$. For a standard $12$-layer ViT and $K=4$ tasks, this results in a $48$-dimensional search space. As we will demonstrate, this high-dimensional search space is highly vulnerable to two distinct points of failure under few-shot validation tuning: (1) \emph{optimization error}, since standard local search algorithms struggle to scale to 48 parameters under tight computation budgets, and (2) \emph{generalization error}, as high-dimensional parameters overfit severely to validation noise when data is scarce.

\subsection{Offline Few-Shot Validation Tuning (OFS-Tune)}
Let $f(x; W_{merged}(\theta))$ represent the classification prediction of the merged model on input $x$. Suppose we have access to a tiny, labeled validation dataset $D_{val}$ that is partitioned by task:
\begin{equation}
D_{val} = \bigcup_{k=1}^K D_{val}^k
\end{equation}
Here, each $D_{val}^k$ is a few-shot dataset containing exactly $M$ labeled samples from task $k$ (where $M \in \{5, 10, 20, 50\}$).

\paragraph{Supervised Few-Shot vs. Unsupervised Zero-Shot TTA}
Before presenting the optimization objective, we must address an important methodological and conceptual distinction. Online TTA model merging operates under a \emph{fully unsupervised zero-shot target regime}, requiring absolutely zero labeled target-task samples. In contrast, OFS-Tune operates under a \emph{supervised few-shot regime}. This represents an ``apples-to-oranges'' problem setup comparison.
As \emph{The Methodologist}, we emphasize that this is a pragmatic, real-world trade-off rather than a direct mathematical replacement. While there exist highly restrictive deployment scenarios (e.g., zero-shot domain generalization with strict privacy or proprietary constraints where target labeled data is strictly impossible to acquire), in the vast majority of practical software engineering applications, practitioners can easily procure a tiny labeled validation set (e.g., 5--10 validation samples per task).
Therefore, we argue that OFS-Tune should be established as a mandatory baseline that unsupervised TTA methods must outperform whenever few-shot validation data is available, serving as a robust, simple, and computationally trivial alternative to active online TTA.

\paragraph{Validation Domain Shift and Selection Bias}
In real-world applications, a prominent methodological threat to few-shot validation tuning is \emph{selection bias} or \emph{validation domain shift}, which occurs when the tiny validation dataset $D_{val}^k$ is not perfectly representative of the true, unobserved target test distribution. We formalize this domain shift by introducing a systematic target bias vector $v_{bias} \sim \mathcal{N}(0, \sigma^2_{bias} I)$ to the validation target profiles during offline optimization. This represents a systematic mismatch between the validation set and the test set. We evaluate optimization across various search spaces under increasingly severe systematic shifts (from $\sigma_{bias} = 0\%$ up to $30\%$). This allows us to empirically test whether low-dimensional search spaces act as robust filters against validation selection bias in addition to sample noise.

The validation objective is to find the merging parameters $\theta^*$ that minimize the joint cross-entropy loss over $D_{val}$:
\begin{equation}
\begin{split}
\mathcal{L}_{val}(\theta) = \frac{1}{K} \sum_{k=1}^K & \frac{1}{|D_{val}^k|} \\
& \cdot \sum_{(x, y) \in D_{val}^k} \mathcal{L}_{CE}\left(f(x; W_{merged}(\theta)), y\right)
\end{split}
\end{equation}

Since the validation loss landscape is non-convex and non-differentiable due to the discrete nature of the few-shot sample, gradient-based methods are prone to getting stuck in poor local minima or fitting to transductive noise. However, because our parameterizations (GT-Merge and Poly-Val-Merge) have exceptionally low dimensionality (e.g., $K=4$ or $K(d+1)=8$), we can optimize $\mathcal{L}_{val}(\theta)$ extremely efficiently using derivative-free black-box optimization. Specifically, we employ Scipy's \textbf{Nelder-Mead} simplex algorithm, initialized from the standard uniform baseline ($\alpha_k = 0.3$).

Crucially, our low-dimensional parameterizations provide a \emph{Dual Optimization-Generalization Benefit}: they make the offline optimization problem mathematically tractable for simple local search algorithms while simultaneously acting as analytical low-pass filters that reject sample noise and prevent generalization overfitting. In contrast, running Nelder-Mead on the 48-dimensional unconstrained space suffers from severe optimization failures, leading to poor convergence.

Furthermore, we analyze the scalability of OFS-Tune with respect to the number of tasks $K$. As $K$ scales, derivative-free simplex search algorithms like Nelder-Mead suffer from catastrophic dimensionality bottlenecks. To enable large-scale model merging, we extend OFS-Tune to use differentiable validation optimization. Since the analytical validation loss is differentiable, we can compute exact gradients and employ gradient-based optimizers like \textbf{PyTorch Adam} directly on the weight-space parameters, allowing OFS-Tune to scale smoothly to dozens or hundreds of tasks.

\begin{algorithm}[tb]
\caption{Offline Few-Shot Validation Tuning (OFS-Tune)}
\label{alg:ofs_tune}
\begin{algorithmic}[1]
\STATE {\bfseries Input:} Labeled few-shot validation set $D_{val} = \bigcup_k D_{val}^k$ with $M$ samples per task, expert task vectors $\{V_k\}_{k=1}^K$, base model $W_{base}$.
\STATE {\bfseries Initialize:} Set parameter $\theta^{(0)}$ to represent the uniform baseline.
\STATE {\bfseries Optimize:} 
\STATE $\theta^* = \text{Nelder-Mead}\left(\mathcal{L}_{val}(\theta), \theta^{(0)}\right)$
\STATE {\bfseries Construct Static Merged Model:}
\STATE $W_{merged}^{(l)} = W_{base}^{(l)} + \sum_k \alpha_k(l; \theta^*) V_k^{(l)} \quad \forall l$
\STATE {\bfseries Output:} Static merged weights $W_{merged}$.
\end{algorithmic}
\end{algorithm}

\subsection{Online Test-Time Adaptation (TTA) Baselines}
To understand the limitations of the current paradigm, we briefly outline the formulation of online TTA model merging. During deployment, an unlabeled stream of test data $\{x_t\}_{t=1}^T$ arrives sequentially. Because labels are unavailable, TTA methods optimize the unconstrained layer-wise parameters $\theta = \{\alpha_{k, l}\}$ online by minimizing the entropy of the predictions on the local batch $B_t$:
\begin{equation}
\mathcal{L}_{ent}(\theta) = -\frac{1}{|B_t|} \sum_{x \in B_t} \sum_{c} p_c(x) \log p_c(x)
\end{equation}
where $p_c(x)$ is the predicted probability of class $c$ for input $x$.

Typically, $\theta$ is updated on-the-fly using Adam over multiple steps per batch. We argue that this formulation is highly fragile. First, entropy minimization on unlabeled streams assumes an i.i.d.\ stream; under extreme class imbalance or bursty task streams, the local entropy objective is highly distorted. Second, optimizing a $48$-dimensional parameter space on small unlabeled batches ($|B_t| \le 2$) introduces severe gradient noise, leading to transductive overfitting and catastrophic representational collapse. OFS-Tune avoids these pitfalls completely by optimizing offline, yielding a static, robust, and zero-compute deployment.

\subsection{Controlled Simulation Calibration}
To conduct a rigorous, multi-seed comparative study of model-merging optimization dynamics, we employ a continuous weight-merging simulation landscape calibrated on empirical Vision Transformer (\textbf{ViT-B/32}) classification statistics across MNIST, FashionMNIST, CIFAR-10, and SVHN datasets. The landscape represents a mathematically closed, reproducible testbed modeled on a coupled Model II sensitivity surface, allowing us to cleanly isolate variables like validation sample size, target stream corruptions, optimization failures, and generalization noise. 

As \textbf{The Methodologist}, we maintain absolute intellectual honesty and rigor regarding this abstraction. A comprehensive, formal defense of our controlled simulation framework---including the mathematical necessity of our non-convex cosine entropy surrogate, the realism of deployment stream noise, and our formulation modeling domain diversity and task interference---is presented in Section \ref{appendix:calibration} of the Appendix.
